\documentclass[a4paper,12pt]{article}
\usepackage[T2A]{fontenc}
\usepackage[utf8]{inputenc}
\usepackage[russian]{babel}
\usepackage{amsmath,amssymb,mathtools}
\usepackage{geometry}
\geometry{left=24mm,right=24mm,top=23mm,bottom=23mm}
\usepackage{microtype}
\usepackage{tikz}
\usetikzlibrary{arrows.meta,positioning,shapes.geometric,calc}
\usepackage{tabularx,booktabs,enumitem}
\usepackage{fancyhdr}
\usepackage{parskip}
\usepackage{xcolor}
\usepackage{hyperref}
\definecolor{accent}{HTML}{008080}
\definecolor{darkaccent}{HTML}{005A5A}
\definecolor{textgray}{HTML}{333333}
\definecolor{muted}{HTML}{777777}
\definecolor{slatefill}{HTML}{E8F0F0}
\definecolor{ochrefill}{HTML}{F8F1DE}
\definecolor{rosefill}{HTML}{F3E5E5}
\definecolor{violetfill}{HTML}{EDE9F5}
\color{textgray}
\hypersetup{colorlinks=true,linkcolor=darkaccent,urlcolor=darkaccent,pdftitle={Чек-лист: составные фигуры и множества решений},pdfauthor={Лёвин Артём Александрович}}
\setlength{\parindent}{0pt}
\setlength{\parskip}{6pt plus 1pt}
\setlist[itemize]{leftmargin=5mm,itemsep=3pt,topsep=2pt}
\setlist[enumerate]{leftmargin=6mm,itemsep=7pt,topsep=3pt}
\pagestyle{fancy}
\fancyhf{}
\fancyhead[L]{\footnotesize Геометрия и множества решений}
\fancyhead[R]{\footnotesize Екатерина Скелли}
\fancyfoot[C]{\footnotesize\thepage}
\renewcommand{\headrulewidth}{0.35pt}
\renewcommand{\footrulewidth}{0pt}
\newcommand{\sectionline}{\vspace{1mm}\par\noindent\textcolor{accent}{\rule{\linewidth}{0.45pt}}\par\vspace{1mm}}
\newcounter{unit}
\newcommand{\topic}[1]{\refstepcounter{unit}\par\medskip\noindent{\large\color{darkaccent}\arabic{unit}. #1}\par\sectionline}
\newcommand{\tagline}[1]{\textcolor{darkaccent}{\textbf{#1}}}
\renewcommand{\thesection}{\arabic{section}}
\begin{document}
\thispagestyle{empty}
\begin{center}
\vspace*{30mm}
{\large\textcolor{darkaccent}{МАТЕМАТИКА}}\\[20mm]
{\LARGE Чек-лист}\\[12mm]
{\Large\textcolor{accent}{\begin{tabular}{c}Прямоугольник из квадратов\\Системы и совокупности\end{tabular}}}\\[15mm]
{\normalsize Геометрическое рассуждение и множества решений неравенств}\\[43mm]
{\large Екатерина Скелли}\\[13mm]
13 ноября 2026 г.\\[25mm]
{\small Лёвин Артём Александрович}\\[2mm]
{\small эксклюзивно для Екатерины Скелли}
\vfill
\begin{tikzpicture}[x=1cm,y=1cm]
\draw[accent,line width=0.8pt] (0,0) .. controls (1.35,0.60) and (2.65,-0.35) .. (4.0,0)
  .. controls (5.35,0.48) and (6.10,-0.38) .. (7.5,0);
\end{tikzpicture}
\end{center}
\clearpage

\topic{Главная идея геометрической задачи}
Если прямоугольник составлен из нескольких квадратов, обычно сначала неизвестны почти все их стороны. Выберите \tagline{одну опорную сторону $x$} и последовательно выражайте через неё остальные длины.

\tagline{Три полезных свойства:}
\begin{itemize}
 \item у квадрата все стороны равны;
 \item отрезок можно представить как сумму соседних частей или как разность двух известных длин;
 \item у прямоугольника \textbf{противоположные} стороны равны: $AB=CD$ и $AD=BC$.
\end{itemize}

\textcolor{darkaccent}{\textbf{Модель из занятия.}} Прямоугольник $ABCD$ разбит на девять квадратов разных размеров. Сторона самого маленького квадрата равна $1$. Сторону центрального квадрата обозначаем $x$ (на доске это $E_2E_5$).

\begin{center}
\begin{tikzpicture}[scale=0.215, every node/.style={font=\scriptsize}]
  % Coordinates are the exact 33 x 32 dissection when x=7.
  \fill[violetfill] (0,17) rectangle (15,32);
  \fill[ochrefill] (15,14) rectangle (33,32);
  \fill[slatefill] (0,9) rectangle (8,17);
  \fill[rosefill] (8,10) rectangle (15,17);
  \fill[ochrefill] (0,0) rectangle (9,9);
  \fill[slatefill] (9,0) rectangle (19,10);
  \fill[rosefill] (19,0) rectangle (33,14);
  \fill[violetfill] (15,10) rectangle (19,14);
  \fill[white] (8,9) rectangle (9,10);
  \draw[line width=0.65pt] (0,0) rectangle (33,32);
  \draw[line width=0.45pt] (0,17)--(15,17)--(15,32);
  \draw[line width=0.45pt] (15,14)--(33,14) (19,0)--(19,14) (15,10)--(19,10);
  \draw[line width=0.45pt] (0,9)--(9,9)--(9,0) (8,9)--(8,17);
  \draw[line width=0.45pt] (8,10)--(19,10) (15,10)--(15,17);
  \node[scale=1.1] at (7.5,25) {$2x+1$};
  \node[scale=1.1] at (24,23) {$x+11$};
  \node at (4,13) {$x+1$};
  \node at (11.5,13.5) {$x$};
  \node at (4.4,4.3) {$x+2$};
  \node at (14,5) {$x+3$};
  \node at (26,7) {$x+7$};
  \node[scale=0.9] at (17,12) {$4$};
  \node[anchor=east] at (7.5,8.0) {$1\,\nearrow$};
  \node[anchor=south east] at (0,32) {$A$};
  \node[anchor=south west] at (33,32) {$B$};
  \node[anchor=north west] at (33,0) {$C$};
  \node[anchor=north east] at (0,0) {$D$};
\end{tikzpicture}
\end{center}
\vspace{-1mm}
{\small\textcolor{muted}{Схема сохраняет реальные пропорции найденного разбиения. Числа внутри квадратов --- выражения для длин их сторон.}}

\tagline{Как появляется первое выражение.} Сторона соседнего квадрата состоит из отрезка длины $x$ и единичного отрезка:
\[
 x+1.
\]
После этого следующая сторона может оказаться равной $(x+1)+1=x+2$. Все дальнейшие выражения получаются из таких же равенств и разбиений отрезков.
\clearpage

\topic{Как последовательно найти стороны}
Выберите опорную сторону $x$. При чтении схемы каждую новую длину проверяйте по уже найденным квадратам.
\begin{center}
\begin{tabularx}{0.93\linewidth}{@{}X>{\raggedleft\arraybackslash}p{47mm}@{}}
\toprule
Как получается длина & Длина стороны\\
\midrule
Самый маленький квадрат & $1$\\
Опорный центральный квадрат & $x$\\
Прибавили единичный отрезок & $x+1$\\
Ещё один единичный отрезок & $x+2$\\
Сложили соседние стороны $x$ и $x+1$ & $2x+1$\\
К стороне $x+2$ прибавили $1$ & $x+3$\\
Вычли $x-1$ из $x+3$ & $(x+3)-(x-1)=4$\\
К стороне $x+3$ прибавили $4$ & $x+7$\\
К стороне $x+7$ прибавили $4$ & $x+11$\\
\bottomrule
\end{tabularx}
\end{center}

\tagline{Самый важный шаг --- два выражения для одной длины.}
Внешняя фигура --- прямоугольник. Составим его \textbf{левую} и \textbf{правую} стороны из известных частей:
\begin{align*}
 AD&=(2x+1)+(x+1)+(x+2)=4x+4,\\
 BC&=(x+11)+(x+7)=2x+18.
\end{align*}
Эти стороны противоположные, поэтому $AD=BC$:
\[
4x+4=2x+18\quad\Longrightarrow\quad 2x=14\quad\Longrightarrow\quad\boxed{x=7}.
\]

\tagline{Искомые размеры прямоугольника:}
\begin{align*}
AB&=(2x+1)+(x+11)=3x+12=\boxed{33},\\
BC&=(x+11)+(x+7)=2x+18=\boxed{32}.
\end{align*}
Проверка: $AD=4\cdot7+4=32=BC$. Площадь прямоугольника $33\cdot32=1056$. Сумма площадей всех девяти квадратов также равна $1056$:
\[
1^2+4^2+7^2+8^2+9^2+10^2+14^2+15^2+18^2=1056.
\]
\textcolor{darkaccent}{\textbf{Вывод.}} Мы приравняли длины одних и тех же по величине отрезков, выраженные через общую переменную. Именно это превратило геометрическую задачу в обычное уравнение.
\clearpage

\begin{center}
{\Large\textcolor{darkaccent}{Алгоритм решения задачи о составной фигуре}}\\[6mm]
{\small Применяйте этот маршрут к другим рисункам, составленным из квадратов.}
\end{center}
\vspace{4mm}
\begin{center}
\begin{tikzpicture}[
  >=Latex,
  node distance=9mm,
  proc/.style={draw=darkaccent,rounded corners=2pt,align=center,text width=104mm,minimum height=13mm,inner sep=4pt,font=\small},
  decision/.style={draw=darkaccent,diamond,aspect=3.0,align=center,text width=61mm,inner sep=2.5pt,font=\small},
  arr/.style={->,draw=accent,line width=0.8pt}
]
\node[proc] (start) {Выберите опорную сторону $x$\\у одного из квадратов.};
\node[proc,below=of start] (parts) {Рассматривайте соседние квадраты:\\стороны равны, отрезки складываются или вычитаются.};
\node[proc,below=of parts] (calc) {Выразите стороны доступных квадратов через $x$.};
\node[decision,below=12mm of calc] (gate) {Есть два выражения\\для равных отрезков?};
\node[proc,below=13mm of gate] (eq) {Запишите равенство, решите уравнение и найдите $x$.};
\node[proc,below=of eq] (ans) {Вычислите запрошенные длины.\\Проверьте равенство противоположных сторон.};
\node[proc,right=16mm of gate,text width=37mm,minimum height=20mm] (more) {Найдите следующий доступный квадрат и выразите его сторону.};
\draw[arr] (start)--(parts);
\draw[arr] (parts)--(calc);
\draw[arr] (calc)--(gate);
\draw[arr] (gate)--node[left,font=\scriptsize]{да}(eq);
\draw[arr] (eq)--(ans);
\draw[arr] (gate.east)--node[above,font=\scriptsize]{нет}(more.west);
\draw[arr] (more.north)|-([yshift=4.5mm]calc.east)--(calc.east);
\end{tikzpicture}
\end{center}
\vfill
\begin{center}
\textcolor{darkaccent}{\textbf{Важно:}} у прямоугольника равны противоположные стороны.\\
Соседние стороны в общем случае имеют разную длину.
\end{center}
\clearpage

\topic{Система неравенств: пересечение}
Числовая прямая помогает увидеть множество чисел, удовлетворяющих неравенству.

\tagline{Система} означает, что \textbf{все условия выполняются одновременно}. На прямой это общая часть отмеченных областей.
\[
\begin{cases}
 x\ge -5,\\ x\le 10
\end{cases}
\quad\Longleftrightarrow\quad -5\le x\le 10
\quad\Longleftrightarrow\quad x\in[-5;10].
\]
\begin{center}
\begin{tikzpicture}[x=0.32cm,y=0.65cm]
\draw[-{Latex[length=2mm]},thin] (-9,0)--(14,0) node[right] {$x$};
\foreach \v in {-5,0,10}{\draw (\v,0.13)--(\v,-0.13) node[below=2pt,font=\scriptsize] {$\v$};}
\draw[accent,line width=2.2pt] (-5,0)--(10,0);
\fill[accent] (-5,0) circle (0.17cm);
\fill[accent] (10,0) circle (0.17cm);
\end{tikzpicture}
\end{center}

\topic{Совокупность неравенств: объединение}
\tagline{Совокупность} означает, что достаточно выполнить \textbf{хотя бы одно} из условий. На прямой нужно отметить все подходящие области.
\[
\left[\begin{array}{l}
 x\le -5,\\
 x\ge 10
\end{array}\right.
\quad\Longleftrightarrow\quad
x\in(-\infty;-5]\cup[10;+\infty).
\]
\begin{center}
\begin{tikzpicture}[x=0.32cm,y=0.65cm]
\draw[-{Latex[length=2mm]},thin] (-9,0)--(14,0) node[right] {$x$};
\foreach \v in {-5,0,10}{\draw (\v,0.13)--(\v,-0.13) node[below=2pt,font=\scriptsize] {$\v$};}
\draw[accent,line width=2.2pt,-{Latex[length=2.5mm]}] (-5,0)--(-8.5,0);
\draw[accent,line width=2.2pt,-{Latex[length=2.5mm]}] (10,0)--(13.5,0);
\fill[accent] (-5,0) circle (0.17cm);
\fill[accent] (10,0) circle (0.17cm);
\end{tikzpicture}
\end{center}

\textcolor{darkaccent}{\textbf{Важное уточнение.}} Области в совокупности могут и пересекаться. Если они пересекаются, объедините их в один более длинный промежуток.

\tagline{Как отличить знаки на концах?}
\begin{itemize}
\item $\le$ и $\ge$: точка входит в ответ, круг закрашен, квадратная скобка $[\ ]$.
\item $<$ и $>$: точка исключена, круг пустой, круглая скобка $(\ )$.
\end{itemize}
\clearpage

\topic{Смешанная запись: системы внутри совокупности}
Чтобы перечислить несколько отрезков или лучей, используйте общую совокупность. Каждый ограниченный с двух сторон промежуток задавайте своей системой.
\[
(-\infty;-5]\cup[0;1]\cup[10;+\infty)
\]
можно записать следующим образом:
\[
\left[\begin{array}{l}
 x\le-5,\\[2pt]
 \begin{cases}x\ge0,\\x\le1,\end{cases}\\[5pt]
 x\ge10.
\end{array}\right.
\]
Здесь выполняется один из трёх вариантов. \textbf{Внутри среднего варианта} оба ограничения выполняются одновременно.

\tagline{Обратная задача.} По записи
\[
\left[\begin{array}{l}
 x\le -5,\\[2pt]
 \begin{cases}x\ge3,\\x\le5,\end{cases}\\[5pt]
 \begin{cases}x\ge7,\\x\le10,\end{cases}\\[5pt]
 x\ge15
\end{array}\right.
\]
получаем четыре отдельные области:
\[
(-\infty;-5]\cup[3;5]\cup[7;10]\cup[15;+\infty).
\]

\topic{Типичные ошибки и самопроверка}
\begin{itemize}
\item \textbf{Сравнили соседние стороны прямоугольника.} Проверяйте, что приравниваете именно противоположные.
\item \textbf{Потеряли единичный отрезок.} Выписывайте каждое разложение вида «часть + часть» непосредственно с рисунка.
\item \textbf{Поменяли местами систему и совокупность.} Система --- \emph{все одновременно}; совокупность --- \emph{хотя бы одно}.
\item \textbf{Забыли про граничную точку.} Перед записью ответа посмотрите на строгий или нестрогий знак.
\item \textbf{Не объединили пересекающиеся области.} После решения обязательно упростите итоговое множество.
\end{itemize}
\textcolor{darkaccent}{\textbf{Мини-проверка:}} можете ли Вы объяснить, почему $AD=BC$ в геометрической задаче и почему $[0;1]$ в совокупности задаётся системой из двух неравенств?
\clearpage

\begin{center}
{\Large\textcolor{darkaccent}{Тренировочный блок}}\\[2mm]
{\small Путь А --- 10 задач. От коротких преобразований к составным ситуациям.}
\end{center}
\sectionline
\textcolor{darkaccent}{\textbf{Геометрия и отрезки}}
\begin{enumerate}[label=\textbf{\arabic*.}]
\item Отрезок длиной $x$ продолжили на $1$ единицу. Запишите общую длину.
\item Отрезок длиной $x+3$ состоит из частей $x-1$ и $y$, где $x>1$. Найдите $y$.
\item В прямоугольнике противоположные стороны имеют длины $4x+4$ и $2x+18$. Найдите $x$ и длину этих сторон.
\item В раскладке из девяти квадратов со схемы найдите длины сторон всех квадратов при $x=7$. Выпишите их по возрастанию.
\item Используя полученные в задаче 4 девять длин, найдите сумму площадей квадратов. Сравните её с площадью прямоугольника размером $33\times32$.
\end{enumerate}
\vspace{4mm}
\textcolor{darkaccent}{\textbf{Системы, совокупности и числовая прямая}}
\begin{enumerate}[label=\textbf{\arabic*.},start=6]
\item Представьте промежуток $[-3;8)$ в виде системы неравенств.
\item Решите систему и запишите ответ промежутком:
\[\begin{cases}x\ge-2,\\x<5.\end{cases}\]
\item Решите совокупность и изобразите ответ на оси:
\[\left[\begin{array}{l}x\le-4,\\x>3.\end{array}\right.\]
\item Запишите с помощью одной совокупности (при необходимости используя внутри системы) множество
\[[-6;-1]\cup(2;4]\cup[7;+\infty).\]
\item Решите совокупность, \textbf{объединив пересекающиеся промежутки}:
\[\left[\begin{array}{l}
\begin{cases}x\ge-4,\\x\le1,\end{cases}\\[4pt]
\begin{cases}x\ge0,\\x<5,\end{cases}\\[4pt]
x\ge9.
\end{array}\right.\]
\end{enumerate}
\clearpage

\begin{center}{\Large\textcolor{darkaccent}{Ответы к тренировочному блоку}}\end{center}
\sectionline
\renewcommand{\arraystretch}{1.35}
\begin{center}
\begin{tabularx}{\linewidth}{@{}>{\centering\arraybackslash}p{12mm}X@{}}
\toprule
\textbf{№}&\textbf{Ответ}\\
\midrule
1 & $x+1$.\\
2 & $y=4$.\\
3 & $x=7$, противоположные стороны равны $32$.\\
4 & $1;\ 4;\ 7;\ 8;\ 9;\ 10;\ 14;\ 15;\ 18$.\\
5 & $1056$ квадратных единиц; площади совпадают: $33\cdot32=1056$.\\
6 & $\begin{cases}x\ge-3,\\x<8.\end{cases}$\\
7 & $[-2;5)$.\\
8 & $(-\infty;-4]\cup(3;+\infty)$.\\
9 & $\displaystyle\left[\begin{array}{l}
\begin{cases}x\ge-6,\\x\le-1,\end{cases}\\[3pt]
\begin{cases}x>2,\\x\le4,\end{cases}\\[3pt]
x\ge7.
\end{array}\right.$\\
10 & $[-4;5)\cup[9;+\infty)$.\\
\bottomrule
\end{tabularx}
\end{center}
\vfill
\begin{center}
\textcolor{darkaccent}{\textbf{Итоговая памятка}}\\[4mm]
\begin{tabularx}{0.94\linewidth}{@{}p{51mm}X@{}}
\toprule
\textbf{Сигнал} & \textbf{Полезное действие}\\
\midrule
Много соседних квадратов & Выразите все стороны через одну переменную.\\
Две противоположные стороны & Запишите равенство длин.\\
«И» в условиях & Найдите пересечение: система.\\
«Или» в условиях & Найдите объединение: совокупность.\\
Точка на границе & Проверьте знак и вид скобки.\\
\bottomrule
\end{tabularx}
\par\vspace{12mm}
{\footnotesize Лёвин Артём Александрович --- эксклюзивно для Екатерины Скелли}
\end{center}
\end{document}
